\documentclass[11pt,letterpaper]{article}
\usepackage[T1]{fontenc}
\usepackage[utf8]{inputenc}
\usepackage{lmodern}
\usepackage[margin=1in,headheight=15pt]{geometry}
\usepackage{amsmath,amssymb,amsthm,mathtools}
\usepackage{microtype}
\usepackage{booktabs,array,longtable}
\usepackage[dvipsnames]{xcolor}
\usepackage{enumitem}
\usepackage{fancyhdr}
\usepackage{hyperref}
\usepackage{listings}
\hypersetup{colorlinks=true,linkcolor=MidnightBlue,citecolor=MidnightBlue,
 urlcolor=MidnightBlue,pdftitle={Jacobsthal Products All at Once},
 pdfauthor={Research manuscript prepared with ChatGPT},
 pdfsubject={An order-theoretic construction for Altman's Question 2.16}}
\pagestyle{fancy}
\fancyhf{}
\fancyhead[L]{\small Jacobsthal products all at once}
\fancyhead[R]{\small Research manuscript}
\fancyfoot[C]{\thepage}
\setlength{\parskip}{4pt}
\setlength{\parindent}{1.3em}
\setcounter{tocdepth}{2}
\allowdisplaybreaks[2]
\newtheorem{theorem}{Theorem}[section]
\newtheorem{lemma}[theorem]{Lemma}
\newtheorem{proposition}[theorem]{Proposition}
\newtheorem{corollary}[theorem]{Corollary}
\theoremstyle{definition}
\newtheorem{definition}[theorem]{Definition}
\newtheorem{example}[theorem]{Example}
\theoremstyle{remark}
\newtheorem{remark}[theorem]{Remark}
\newcommand{\J}{\mathbin{\times_{\!J}}}
\newcommand{\natsum}{\mathbin{\oplus}}
\newcommand{\natprod}{\mathbin{\otimes}}
\newcommand{\FF}{\mathfrak F}
\newcommand{\LL}{\mathfrak L}
\newcommand{\BB}{\mathfrak B}
\newcommand{\CC}{\mathfrak C}
\newcommand{\otp}{\operatorname{otp}}
\newcommand{\CNF}{\operatorname{CNF}}
\newcommand{\supp}{\operatorname{supp}}
\newcommand{\fs}{\operatorname{FS}}
\newcommand{\colex}{\operatorname{Colex}}
\newcommand{\lead}{\operatorname{lm}}
\newcommand{\one}{\mathbf 1}
\newcommand{\tail}{\operatorname{tail}}
\newcommand{\concat}{\mathbin{{}^{\frown}}}
\newcommand{\emptyword}{\langle\rangle}
\newcommand{\ep}{\varepsilon_0}
\lstset{basicstyle=\ttfamily\small,columns=fullflexible,breaklines=true,
 frame=single,rulecolor=\color{black!25},showstringspaces=false}

\begin{document}
\begin{titlepage}
\centering
\vspace*{0.7in}
{\large\sffamily ORDER THEORY \quad / \quad ORDINAL ARITHMETIC\par}
\vspace{0.35in}
{\Huge\bfseries Jacobsthal Products\\[5pt] All at Once\par}
\vspace{0.22in}
{\Large A Cantor-block construction and a finite-tail reduction\par}
\vspace{0.28in}
{\large A research attempt on Altman's Question 2.16\par}
\vspace{0.22in}
{20 September 2026\par}
\vspace{0.32in}
\begin{minipage}{0.92\textwidth}
\begin{abstract}
Jacobsthal multiplication is obtained by transfinitely iterating natural
addition. Its associativity is classical, but Altman asked for a direct
interpretation of products of several, and even transfinitely many, factors
without defining the answer by parenthesized iteration or a limit. We give a
Cantor-normal-form-based construction on the actual Cartesian product in the
finite case and on the finite-support Cartesian product in the transfinite
case.

For finitely many factors, a tuple is assigned a pivot: the last coordinate,
other than the first, lying in a positive-exponent Cantor block. The resulting
finitely many cells receive ordinary reverse-lexicographic orders and are
placed in decreasing order of their monomial types. A block comparison proves
compatibility with finite regrouping without assuming Jacobsthal
associativity. The construction is a componentwise linear extension and is
maximal among orders subject to its within-cell constraints.

For arbitrary nonzero factors, ordinary and Jacobsthal partial products have
the same Cantor exponents and the same leading monomial. A uniform bound
$O_\xi\leq J_\xi<O_\xi\cdot2$ implies equality at every limit with cofinally
many nonunit factors. Removing units leaves an active index of the form
$\lambda+n$, where $\lambda$ is zero or a limit and $n$ is finite. The entire
product is an ordinary finite-support prefix followed by a single finite
pivot product. We also obtain an exact three-state test for equality with the
ordinary product and an order-theoretic proof of the binary natural-product
upper bound. Exact computational checks accompany, but do not replace, the
proofs.
\end{abstract}
\vspace{0.12in}
\noindent\textbf{Status.} This is an unrefereed, AI-assisted research
manuscript, not a priority certification or a proof-assistant formalization.
The proposed answer uses a precise CNF-based meaning of ``all at once.''
The scope of the claim and the limitations of the literature search are
stated explicitly in Sections~\ref{sec:problem} and~\ref{sec:scope}.
\end{minipage}
\end{titlepage}

\tableofcontents
\newpage

\section{The selected problem and the proposed answer}\label{sec:problem}

\subsection{A published question, rather than a newly invented conjecture}

The selected problem is Question~2.16 in Harry Altman's
\emph{Intermediate arithmetic operations on ordinal numbers}
\cite{altman}. In the notation of this manuscript, it asks for an
order-theoretic meaning of $a\J b\J c$ that does not first choose parentheses,
and an interpretation of a transfinitely indexed Jacobsthal product other
than its recursive limit definition. The operation $\J$ is Altman's
$\times$; we reserve ordinary juxtaposition for ordinary ordinal
multiplication.

The associativity identity itself is not an open problem. It is a classical
result of Jacobsthal, reproduced in Altman's Theorem~2.13. Nor is the binary
Cantor-normal-form formula new. The objective here is to construct the
underlying ordered objects directly and explain how the already known
arithmetic emerges from them.

There is relevant prior work on infinitary natural sums and products.
Lipparini's work on transfinite natural sums provides related
order-theoretic descriptions \cite{lip-sums}; his countable natural-product
construction combines finite natural-product behavior with ordinary
finite-support products \cite{lip-product}. The ordinary finite-support
model used below is standard, not a new construction. A recent general
framework for noncommutative infinitary operations appears in
\cite{lip-semigroup}. Our proposed construction concerns Jacobsthal
multiplication specifically.

A targeted literature search on 20 September 2026 located the published
question and these related works, but did not locate a subsequent solution
to Question~2.16. This is a statement about the search, not a proof of
worldwide novelty. The mathematical arguments below can be assessed
independently of that unresolved bibliographic question.

\subsection{What ``all at once'' means here}

We impose the following concrete requirements. The object must be defined
on tuples, or on finite-support functions when the index is transfinite.
Its comparison rule must use only the input well-orders and explicitly
specified ordinary order constructions. Neither a parenthesized
Jacobsthal product nor the supremum of its partial values may occur in the
definition of that comparison rule. We also require invariance under order
isomorphisms of the inputs and extension of the componentwise partial
order.

These requirements permit the canonical Cantor decomposition of each input
well-order. They do not require the comparison rule to be preserved by
arbitrary order embeddings. They also do not require the usual flattening
bijection between nested tuples to be the associativity isomorphism. These
distinctions are substantive, not merely terminological.

\subsection{The shape of the answer}

For a finite list, we construct a well-order
$\FF(A_1,\ldots,A_r)$ on $A_1\times\cdots\times A_r$.
Write $F(a_1,\ldots,a_r)$ for its order type. Its definition, given in
Section~\ref{sec:finite}, is a single rule for all finite arities.
Section~\ref{sec:grouping} proves
\begin{equation}\label{eq:intro-group}
 \FF(U\concat V)\cong \FF\bigl(\FF(U),\FF(V)\bigr),
\end{equation}
where $U,V$ are finite lists of well-orders. This proof compares monomial
cells and does not assume associativity of $\J$.

For a sequence $(a_i)_{i<\kappa}$ of nonzero ordinals, remove all factors
$1$ and let the remaining, or \emph{active}, index have order type
$\tau=\lambda+n$, with $\lambda=0$ or $\lambda$ a nonzero limit ordinal,
and $n<\omega$. Denote the reindexed active factors by $(b_\eta)_{\eta<\tau}$.
The main reduction is
\begin{equation}\label{eq:intro-reduction}
 \boxed{\quad
 J((a_i)_{i<\kappa})
   =p\,F(b_\lambda,\ldots,b_{\lambda+n-1}),\qquad
 p=\prod_{\eta<\lambda}b_\eta .\quad}
\end{equation}
Here the product defining $p$ is \emph{ordinary}; $p=1$ when $\lambda=0$.
When $\lambda>0$, $p$ is a pure power of $\omega$.
The right side is realized by a direct order on finite-support functions:
compare the finite terminal tuples using $\FF$; when they agree, compare
the prefix functions by their largest differing coordinate. Thus the
\emph{definition of the final ordered set} contains no limit of Jacobsthal
products.

Two further results make the reduction checkable and useful. At every
partial stage ordinary and Jacobsthal products have identical CNF exponent
supports and identical leading monomials. Moreover, equality of the two
full products is decided by a three-state computation on the finite active
suffix alone. Both statements are proved below for arbitrary ordinals,
not merely for the computable ordinals used in the tests.

\section{Conventions and the binary arithmetic}\label{sec:background}

\subsection{Cantor blocks, tails, and three multiplications}

For a nonzero ordinal $a$, write
\begin{equation}\label{eq:cnf}
 a=\omega^{\rho_1}c_1+\cdots+\omega^{\rho_s}c_s,
 \qquad \rho_1>\cdots>\rho_s,\qquad 0<c_i<\omega.
\end{equation}
Its degree is $d(a)=\rho_1$, its leading monomial is
$\lead(a)=\omega^{\rho_1}c_1$, and its CNF exponent support is
$\supp_{\CNF}(a)=\{\rho_1,\ldots,\rho_s\}$.
A \emph{monomial} has one distinct exponent, with any positive finite
coefficient. A \emph{pure power} has the more restrictive form $\omega^\rho$.
In particular, every positive finite ordinal is a monomial, but only $1$
is a finite pure power.

There is a unique decomposition
\begin{equation}\label{eq:tail}
 a=a^++t(a),\qquad t(a)<\omega,
\end{equation}
where $a^+$ is zero or a nonzero limit ordinal. The number $t(a)$ is the
coefficient of exponent $0$, or zero if that exponent is absent.
For a well-order $A$ of type $a$, let $T_A$ be its final finite interval of
length $t(a)$.

Expand every coefficient in \eqref{eq:cnf} into repeated unit blocks:
\begin{equation}\label{eq:unit-cnf}
 a=\omega^{\sigma_1}+\cdots+\omega^{\sigma_m},
 \qquad \sigma_1\geq\cdots\geq\sigma_m.
\end{equation}
Transporting the corresponding intervals along the unique isomorphism
$A\cong a$ gives canonical convex blocks $B_{A,q}$ of type
$\omega^{\sigma_q}$. Blocks of exponent $0$ are singletons and together
form $T_A$. No choice of a noncanonical presentation is involved.

Natural sum $\oplus$ adds CNF coefficients at equal exponents. Natural
product $\otimes$ multiplies the resulting formal polynomials, using
natural addition for the exponents. In particular,
\begin{equation}\label{eq:finite-natural}
 a\otimes n=\sum_{j=1}^{s}\omega^{\rho_j}(c_jn)
 \quad(n<\omega),
\end{equation}
with zero terms omitted. In contrast, for $n>0$,
\begin{equation}\label{eq:finite-ordinary}
 an=\omega^{\rho_1}(c_1n)+\omega^{\rho_2}c_2+\cdots+
      \omega^{\rho_s}c_s.
\end{equation}
Thus ordinary multiplication by a positive finite integer changes only
the leading coefficient, whereas natural multiplication changes them all.

Jacobsthal multiplication is defined by
\begin{equation}\label{eq:J-def}
 a\J0=0,\qquad a\J(\beta+1)=(a\J\beta)\oplus a,
 \qquad a\J\lambda=\sup_{\beta<\lambda}(a\J\beta)
\end{equation}
at nonzero limit $\lambda$. We use these equations only to identify our
constructed objects with the standard operation.

\subsection{Ordinary absorption and a self-contained binary calculation}

\begin{lemma}\label{lem:ordinary}
Let $a>0$, $d=d(a)$, and $\sigma>0$. Then
\begin{equation}\label{eq:absorb}
 a\omega^\sigma=\omega^{d+\sigma}.
\end{equation}
For $b=b^++t$ as in \eqref{eq:tail},
\begin{equation}\label{eq:ordinary-binary}
 ab=\omega^d b^++at.
\end{equation}
Consequently $d(ab)=d(a)+d(b)$ when $a,b>0$.
\end{lemma}

\begin{proof}
The equality $a\omega=\omega^{d+1}$ follows by taking the supremum of the
finite multiples in \eqref{eq:finite-ordinary}. Induction on positive
$\sigma$, using ordinary associativity at successors and continuity at
limits, gives \eqref{eq:absorb}. Expanding $b^+$ into its positive-exponent
CNF terms and using left distributivity of ordinary multiplication gives
\eqref{eq:ordinary-binary}. Its highest exponent is $d(a)+d(b)$ when
$b$ is infinite. When $b$ is finite and positive, the highest exponent
remains $d(a)=d(a)+0$.
\end{proof}

\begin{proposition}[Classical Jacobsthal formula]\label{prop:binary}
For $a>0$, $d=d(a)$, and $b=b^++t$, one has
\begin{equation}\label{eq:J-binary}
 \boxed{a\J b=\omega^d b^++(a\otimes t).}
\end{equation}
This is the formula in Jacobsthal's theorem, restated in
\cite[Theorem~2.10]{altman}. The proof below is included to isolate precisely
the arithmetic input needed later.
\end{proposition}

\begin{proof}
Let $H(b)$ denote the right side of \eqref{eq:J-binary}. Plainly $H(0)=0$.
All exponents in $\omega^d b^+$ are strictly greater than $d$. Therefore
adding $a$ naturally changes only the finite-tail contribution, and
\[
 H(b)\oplus a=\omega^d b^++(a\otimes(t+1))=H(b+1).
\]
Now let $\lambda$ be a nonzero limit. Equations
\eqref{eq:finite-natural}--\eqref{eq:ordinary-binary} imply $ab\leq H(b)$.
For $b<\lambda$, every exponent in $a\otimes t(b)$ is at most $d$, so
\[
 H(b)<\omega^d b^++\omega^{d+1}
      =\omega^d(b^++\omega)\leq\omega^d\lambda.
\]
The last inequality uses $b^++\omega\leq\lambda$, since $\lambda$ is a
limit larger than $b$. On the other hand, ordinary continuity gives
\[
 \sup_{b<\lambda}ab=a\lambda=\omega^d\lambda.
\]
Hence $\sup_{b<\lambda}H(b)=\omega^d\lambda=H(\lambda)$. The defining
recursion \eqref{eq:J-def} proves the assertion.
\end{proof}

\begin{corollary}\label{cor:binary-equality}
For nonzero $a,b$,
\begin{equation}\label{eq:binary-equality}
 a\J b=ab
 \quad\Longleftrightarrow\quad
 t(b)\leq1\ \text{or $a$ is a monomial}.
\end{equation}
In particular, $a\J\omega^\sigma=a\omega^\sigma$ for $\sigma>0$,
and $\omega^\delta\J b=\omega^\delta b$ for every $b$.
\end{corollary}

\begin{proof}
The positive-exponent contributions in \eqref{eq:ordinary-binary} and
\eqref{eq:J-binary} coincide and lie above all exponents in the remaining
terms. The two remaining terms agree when $t=0$ or $1$. For $t\geq2$,
comparison of \eqref{eq:finite-natural} and \eqref{eq:finite-ordinary}
shows that they agree exactly when there are no lower CNF terms of $a$.
\end{proof}

\begin{remark}
Ordinary addition in exponents is essential. It is strictly increasing
in its right argument but need not be strictly increasing in its left
argument. We repeatedly use $\delta+\sigma>\delta$ for $\sigma>0$.
We never replace the ordinary prefix-degree sums below by natural sums.
All statements here make sense beyond $\ep$, even when an ordinal appears
as its own leading exponent. Only the testing implementation is restricted
to hereditary CNF expressions below $\ep$.
\end{remark}

\section{The finite pivot construction}\label{sec:finite}

\subsection{A partition of the Cartesian product}

Let $A_1,\ldots,A_r$ be nonempty well-orders with types
$a_1,\ldots,a_r$, where $r\geq1$. Write $B_{j,q}$ for the unit CNF blocks
of $A_j$, with exponents $\sigma_{j,q}$, and put $T_j=T_{A_j}$.
Set
\begin{equation}\label{eq:prefix-degree}
 D_0=0,\qquad D_j=d(a_1)+\cdots+d(a_j).
\end{equation}
These are ordinary finite sums.

\begin{definition}[Pivot and cell]\label{def:pivot}
For a tuple $x=(x_1,\ldots,x_r)$, its pivot is the largest index
$j\in\{2,\ldots,r\}$ for which $x_j\notin T_j$, if such an index exists;
otherwise its pivot is $1$.

A cell with pivot $j\geq2$ has the form
\begin{equation}\label{eq:cell}
 C(j,q;u_{j+1},\ldots,u_r)
 =A_1\times\cdots\times A_{j-1}\times B_{j,q}
      \times\{u_{j+1}\}\times\cdots\times\{u_r\},
\end{equation}
where $\sigma_{j,q}>0$ and $u_i\in T_i$. With pivot $1$, the same formula
means $B_{1,q}\times\{u_2\}\times\cdots\times\{u_r\}$, and all blocks
of $A_1$, including exponent-zero blocks, are allowed.
\end{definition}

The cells are disjoint and cover the product. Indeed, the pivot is unique,
its block is unique, and every later coordinate is necessarily in its
finite tail. There are finitely many cells: each factor has finitely many
unit blocks and each $T_i$ is finite. Some possible pivot indices contribute
no cells because a required later tail is empty.

Within a cell use ordinary \emph{colexicographic} order: compare the first
$j$ coordinates at the largest coordinate where they differ. This is the
usual order model of an ordinary ordinal product. Define the exponent of
the cell by
\begin{equation}\label{eq:cell-exponent}
 e(C)=D_{j-1}+\sigma_{j,q}.
\end{equation}
For $j=1$ this is simply $\sigma_{1,q}$.

\begin{lemma}\label{lem:cell-type}
Every cell, with its colexicographic order, has type $\omega^{e(C)}$.
\end{lemma}

\begin{proof}
For $j=1$ this is the definition of a unit block. For $j\geq2$, the ordinary
product of the first $j-1$ factors has degree $D_{j-1}$ by
Lemma~\ref{lem:ordinary}. Its product with $\omega^{\sigma_{j,q}}$ is
$\omega^{D_{j-1}+\sigma_{j,q}}$, because the pivot exponent is positive.
The fixed later coordinates do not affect the order type.
\end{proof}

\begin{lemma}[Separation of pivot layers]\label{lem:pivot-separation}
If cells $C$ and $C'$ have pivots $j<k$, then $e(C)<e(C')$.
For cells with the same pivot, comparison of their exponents $e(C)$ is
exactly comparison of their input block exponents $\sigma_{j,q}$.
\end{lemma}

\begin{proof}
For the first assertion,
\[
 e(C)\leq D_j\leq D_{k-1}<D_{k-1}+\sigma_{k,q'}=e(C'),
\]
because $k\geq2$ and its pivot block has positive exponent. The second
assertion is strict monotonicity of addition by a fixed ordinal on the
left. In particular, cell exponents from different pivot layers never
collide, even when ordinary exponent addition absorbs a prefix.
\end{proof}

\subsection{The comparison rule}

\begin{definition}[Finite all-at-once order]\label{def:F}
Order the cells first by \emph{decreasing pivot index} $j$, then by
\emph{decreasing input block exponent} $\sigma_{j,q}$. For ties, order
by the colexicographic order of the finite tuple
$(u_{j+1},\ldots,u_r)$, then by increasing block index $q$.
Inside each cell retain its colexicographic order. Denote the concatenation
by $\FF(A_1,\ldots,A_r)$ and its type by $F(a_1,\ldots,a_r)$.

For the empty list put $\FF(\emptyword)=\one$, a singleton. If a factor is
empty, put $\FF(A_1,\ldots,A_r)=\varnothing$.
\end{definition}

The order is a well-order because it is a finite concatenation of
well-orders. By Lemma~\ref{lem:pivot-separation}, it also lists the cells
in decreasing order of their monomial exponents $e(C)$. The tie rule is
deterministic; different reasonable tie rules among equal monomial types
need not give the same pointwise order, but give the same order type.
The chosen rule agrees with ordinary colex when all factors are finite.
For a single factor it gives its original well-order.

Most importantly, Definition~\ref{def:F} never evaluates a Jacobsthal
product. It uses canonical convex blocks, finite tags, and ordinary
colexicographic comparison. The prefix-degree sums $D_j$ are needed
only to \emph{prove the resulting order type}, not to compare tuples.
Thus the comparison rule is even more elementary than a rule that first
computes all the monomial types, and is not a recursive definition
disguised as an $r$-ary notation.

\begin{theorem}[Componentwise extension]\label{thm:extension}
The order $\FF(A_1,\ldots,A_r)$ extends the componentwise partial order on
$A_1\times\cdots\times A_r$.
\end{theorem}

\begin{proof}
Suppose $x_i\leq y_i$ for every $i$, and let $j_x,j_y$ be the pivots.
A coordinate in a final finite interval cannot leave that interval by
increasing. Hence $j_y\leq j_x$.

First suppose $j_y<j_x$. Since $j_x\geq2$, its pivot exponent is positive.
The exponent of the cell of $y$ satisfies
\[
 e(C_y)\leq D_{j_y}\leq D_{j_x-1}
          <D_{j_x-1}+\sigma_{j_x,q_x}=e(C_x).
\]
Here the first inequality also covers $j_y=1$. The cell of $x$ is therefore
placed before that of $y$.

Now suppose $j_x=j_y=j$. Increasing $x_j$ cannot decrease its block index,
and so cannot increase its block exponent. Strict monotonicity of ordinary
addition in the right argument gives $e(C_y)\leq e(C_x)$. A strict
inequality again puts $x$ first. In the equality case the pivot exponents
are equal. The later tail tuples are coordinatewise ordered, and hence are
ordered in their colexicographic tie rule. If they agree, the block index
of $x_j$ is no larger than that of $y_j$. Finally, if both tuples lie in
the same cell, its internal colex order respects every coordinatewise
comparison. These cases prove $x\leq_{\FF}y$.
\end{proof}

\subsection{The explicit all-arity formula}

For $j\leq r$, write
$T_{j,r}=\prod_{i=j}^{r}t(a_i)$, an ordinary finite integer product, and
set $T_{r+1,r}=1$. Terms with coefficient zero are omitted.

\begin{theorem}[Cell formula]\label{thm:formula}
For nonzero $a_1,\ldots,a_r$,
\begin{align}\label{eq:pivot-formula}
 F(a_1,\ldots,a_r)
 &=a_1\otimes T_{2,r}\notag\\
 &\quad\oplus
 \bigoplus_{j=2}^{r}
 \ \bigoplus_{\substack{(\sigma,c)\in\CNF(a_j)\\\sigma>0}}
      \omega^{D_{j-1}+\sigma}\bigl(cT_{j+1,r}\bigr).
\end{align}
All natural sums displayed here are finite.
\end{theorem}

\begin{proof}
Every unit block of $a_1$ gives one pivot-$1$ cell for each choice of later
tail coordinates, hence $T_{2,r}$ copies of its monomial. Every
positive-exponent unit block of $a_j$ gives $T_{j+1,r}$ cells of degree
$D_{j-1}+\sigma$. Lemma~\ref{lem:cell-type} identifies the types. Sorting
the finitely many monomials in decreasing exponent produces their natural
sum. Collecting repeated unit blocks yields \eqref{eq:pivot-formula}.
\end{proof}

\begin{corollary}[Binary identification]\label{cor:F-binary}
For all ordinals $a,b$,
\[
 F(a,b)=a\J b.
\]
\end{corollary}

\begin{proof}
Zero factors are immediate. The positive blocks of $b$ contribute
$\omega^{d(a)}b^+$. Their exponents are strictly greater than $d(a)$, so
they precede every term of the contribution $a\otimes t(b)$. Therefore
the cell formula is exactly \eqref{eq:J-binary}.
\end{proof}

\begin{lemma}[Degree and terminal cardinality]\label{lem:metadata}
For a nonempty finite list of nonzero factors,
\begin{equation}\label{eq:metadata}
 d(F(a_1,\ldots,a_r))=D_r,
 \qquad t(F(a_1,\ldots,a_r))=\prod_{i=1}^{r}t(a_i).
\end{equation}
The final finite interval of $\FF$ consists precisely of the tuples all
of whose coordinates belong to their final finite intervals.
\end{lemma}

\begin{proof}
If all factors are finite, the claims follow directly. Otherwise let $h$
be the last index with $d(a_h)>0$. All later factors are positive finite
ordinals, so a cell using the leading block at $h$ exists. Its exponent
is $D_{h-1}+d(a_h)=D_r$. Every other cell has exponent at most $D_r$,
which proves the degree assertion. This argument includes $h=1$.

A pivot at $j\geq2$ always has positive exponent. Thus exponent-zero cells
arise exactly from an exponent-zero block at index $1$ and choices in all
later finite tails. They are singletons, and their number is the stated
product. Every positive-exponent cell precedes them.
\end{proof}

\section{Regrouping and a restricted maximality principle}\label{sec:grouping}

\subsection{Why associativity follows from the construction}

\begin{theorem}[Finite concatenation]\label{thm:grouping}
For finite lists $U,V$ of well-orders,
\begin{equation}\label{eq:grouping}
 \FF(U\concat V)\cong\FF\bigl(\FF(U),\FF(V)\bigr).
\end{equation}
The isomorphism is the unique order isomorphism between the two well-orders.
The proof does not use Jacobsthal associativity.
\end{theorem}

\begin{proof}
Empty factors give empty orders on both sides. An empty list gives a
singleton, which is a unit for the binary construction by its definition
or by \eqref{eq:J-binary}. We may therefore suppose both lists are
nonempty and all factors are nonzero.

Let $\delta_U$ be the ordinary sum of degrees in $U$. By
Lemma~\ref{lem:metadata}, it is the degree of $\FF(U)$. The canonical
unit CNF blocks of $\FF(V)$ are precisely its cells, listed in the order
of their exponents and tags: each such cell has a pure monomial type.
We compare the monomial cells obtained on the right with those on the
left in two groups.

A positive-exponent cell of $\FF(V)$ has some exponent $\eta>0$. In the
outer binary construction, it contributes one cell of type
$\omega^{\delta_U+\eta}$. If its pivot was the $j$th coordinate of $V$,
then
\[
 \eta=D^{V}_{j-1}+\sigma_{j,q},
 \qquad
 \delta_U+\eta=(\delta_U+D^{V}_{j-1})+\sigma_{j,q}.
\]
The second expression is exactly the exponent of the corresponding global
cell whose pivot is in $V$. Its later tail coordinates have not changed.
When the local pivot is $1$, positivity means that its first block has
positive exponent, as required for a global pivot beyond the first
coordinate. Thus this correspondence preserves all multiplicities.

Next consider the exponent-zero cells of $\FF(V)$. By
Lemma~\ref{lem:metadata}, these are its terminal points, in bijection with
tuples of terminal points of the factors in $V$. For each such point, the
outer construction takes every unit block of $\FF(U)$. These are exactly
the cells of $\FF(U)$, and their exponents remain unchanged. They match
the global cells whose pivot is in $U$, with all $V$-coordinates terminal.

We have matched the two finite multisets of pure monomial cell types.
Each construction sorts this multiset in decreasing exponent. Consequently
the ordered runs of equal exponents have the same finite lengths.
Match the first block of each run on one side to the first block on the
other, and so on. Inside matched blocks take the unique isomorphism of
well-orders of type $\omega^\eta$. Their union is the required order
isomorphism.
\end{proof}

\begin{corollary}\label{cor:finite-assoc}
For every finite list, $F(a_1,\ldots,a_r)$ equals every parenthesized
Jacobsthal product of its entries. In particular, $\J$ is associative.
All coherence diagrams made from the order isomorphisms in
Theorem~\ref{thm:grouping} commute.
\end{corollary}

\begin{proof}
Use Corollary~\ref{cor:F-binary} at each binary vertex and
Theorem~\ref{thm:grouping} to identify the result with the common
all-arity object. A well-order has at most one order isomorphism to a
given well-order, so two composites with the same endpoints coincide.
\end{proof}

\begin{remark}[The flattening bijection is not the associator]\label{rem:flatten}
The assertion is an order isomorphism, not equality of comparison rules
under naive tuple flattening. Take factors $\omega+1,2,\omega$. The
direct three-factor construction has a single pivot-$3$ cell, ordered by
ordinary colex. It places $(\omega,0,0)$ before $(0,1,0)$. In the nested
order $\FF(\FF(\omega+1,2),\omega)$, the inner order places $(0,1)$
before $(\omega,0)$, reversing that comparison. Both total orders have
type $\omega^2$; their unique isomorphism is not the flattening identity.
This is precisely why the block-isomorphism argument is needed.
\end{remark}

\subsection{A maximum without building the target value into the definition}

The preceding definition selects an order. There is also an extremal
characterization which removes the significance of its particular tie
rule.

\begin{lemma}[Finite interleaving bound]\label{lem:shuffle}
Let a set be partitioned into finitely many well-orders $C_1,\ldots,C_m$
of types $\alpha_1,\ldots,\alpha_m$. Every linear order on their union
whose restrictions are the given cell orders is a well-order of type at
most $\alpha_1\oplus\cdots\oplus\alpha_m$.
\end{lemma}

\begin{proof}
For a point $z$ and a cell $C_i$, let $r_i(z)$ be the order type of the
points of $C_i$ preceding $z$ in the proposed linear order. This is an
initial segment of $C_i$, so $r_i(z)\leq\alpha_i$. If $z<z'$, all these
ranks are nondecreasing, and the rank of the cell containing $z$ strictly
increases: that cell counts $z$ before $z'$ but not before itself.
Strict monotonicity of finite natural sum therefore makes
\[
 z\longmapsto r_1(z)\oplus\cdots\oplus r_m(z)
\]
strictly increasing. Its value is strictly below
$\alpha_1\oplus\cdots\oplus\alpha_m$, because at $z$ the rank in its
own cell is strictly smaller than the full cell type. Thus it embeds the
linear order into the indicated ordinal, proving both assertions.
\end{proof}

\begin{theorem}[Pivot-cell maximality]\label{thm:maximality}
Among linear extensions of the componentwise order on
$A_1\times\cdots\times A_r$ whose restriction to every pivot cell is
ordinary colexicographic order, $\FF(A_1,\ldots,A_r)$ has the largest
possible order type.
\end{theorem}

\begin{proof}
By Theorem~\ref{thm:extension}, $\FF$ is an admissible order. Every
admissible order is bounded by the finite natural sum of the cell types,
by Lemma~\ref{lem:shuffle}. Since those types are pure powers, concatenating
them in decreasing exponent attains that bound. This is exactly $\FF$.
\end{proof}

The constraints in this theorem refer only to the input CNF blocks and
the ordinary orders inside explicitly described subsets. They do not
refer to $\J$ or to the desired order type. The result is therefore more
than the tautology that a computed ordinal can order a set of suitable
cardinality. It remains a \emph{restricted} maximality theorem: if the
within-cell restrictions are removed, the maximum is generally larger.

\subsection{An order-theoretic comparison with natural multiplication}

\begin{corollary}\label{cor:natural-upper}
For every finite list,
\[
 F(a_1,\ldots,a_r)\leq a_1\otimes\cdots\otimes a_r.
\]
In particular, $a\J b\leq a\otimes b$ has an order-theoretic proof.
\end{corollary}

\begin{proof}
Carruth's theorem identifies the natural product of finitely many
ordinals as the maximum order type of a linear extension of their
componentwise Cartesian product \cite{carruth}; see also the explicit
statement in \cite[Section~3.3]{lip-product}. Our construction is one such
extension, by Theorem~\ref{thm:extension}. Zero factors are immediate.
\end{proof}

The numerical inequality is already known. What this corollary supplies
is the interpretation requested for the \emph{binary multiplication
part} of Altman's Question~5.2 \cite{altman}. It is not a claimed solution
to all the exponentiation and infinitary comparison questions grouped
there.

\section{Synchronization of ordinary and Jacobsthal products}\label{sec:sync}

\subsection{The common-support invariant}

Fix a sequence $(a_i)_{i<\kappa}$ of \emph{nonzero} ordinals. Define
ordinary and Jacobsthal partial products by
\begin{align}\label{eq:partial-products}
 O_0=J_0&=1,\notag\\
 O_{\xi+1}&=O_\xi a_\xi,&
 J_{\xi+1}&=J_\xi\J a_\xi,
\end{align}
and at nonzero limits use the suprema of the preceding partial products.
All factors are at least $1$, so both sequences are nondecreasing.
The treatment of a zero factor is separated in
Section~\ref{subsec:zero}; a supremum formula for nondecreasing products
must not be used carelessly after multiplication by zero.

\begin{lemma}[One-step synchronization]\label{lem:one-step}
Suppose $x,y>0$ have the same CNF exponent support and leading monomial,
and every CNF coefficient of $x$ is at most the corresponding coefficient
of $y$. For every $b>0$, the pair $xb,y\J b$ has the same properties.
If $t(b)=0$, then $xb=y\J b$.
\end{lemma}

\begin{proof}
Let $d=d(x)=d(y)$ and $b=b^++t$. The two outputs are
\[
 \omega^d b^++xt,
 \qquad
 \omega^d b^++(y\otimes t).
\]
When $t=0$ they coincide. When $t>0$, the terms contributed by $b^+$
are identical and have exponents strictly above $d$. In the remaining
terms, $xt$ multiplies the common leading coefficient by $t$ and leaves
all lower coefficients of $x$ unchanged. The ordinal $y\otimes t$
multiplies every coefficient of $y$ by $t$. Thus the remaining exponent
supports are identical and all coefficient inequalities are preserved.
If $b^+\neq0$, the leading monomial comes from the identical upper part.
If $b^+=0$, it comes from the identically scaled old leading monomial.
\end{proof}

\begin{lemma}[A uniform one-step-scale bound]\label{lem:factor-two}
If $x,y>0$ have the same leading monomial and $x\leq y$, then
\[
 x\leq y<x\cdot2.
\]
\end{lemma}

\begin{proof}
Write their leading monomial as $\omega^d c$, with $c\geq1$ finite.
Both ordinals are less than $\omega^d(c+1)$, while
\[
 \omega^d(c+1)\leq\omega^d(2c)\leq x\cdot2.
\]
The first of these inequalities also covers $d=0$. The result follows.
\end{proof}

\begin{theorem}[Synchronization at every ordinal stage]\label{thm:sync}
For every $\xi\leq\kappa$,
\begin{equation}\label{eq:support-sync}
 \supp_{\CNF}(O_\xi)=\supp_{\CNF}(J_\xi),
 \qquad \lead(O_\xi)=\lead(J_\xi).
\end{equation}
Every coefficient of $O_\xi$ is at most its corresponding coefficient in
$J_\xi$, and consequently
\begin{equation}\label{eq:uniform-bound}
 \boxed{O_\xi\leq J_\xi<O_\xi\cdot2.}
\end{equation}
If $\lambda\leq\kappa$ is a nonzero limit and
$\{i<\lambda:a_i>1\}$ is cofinal in $\lambda$, then
\begin{equation}\label{eq:limit-equality}
 O_\lambda=J_\lambda.
\end{equation}
\end{theorem}

\begin{proof}
We prove all the assertions simultaneously by transfinite induction.
They hold at stage $0$. The successor step follows from
Lemma~\ref{lem:one-step}; coefficient domination implies ordinal domination,
and Lemma~\ref{lem:factor-two} gives the strict upper bound.

Let $\lambda$ be a nonzero limit. If the indices $i<\lambda$ with $a_i>1$
are bounded, choose a stage below $\lambda$ after all of them. Every later
factor is $1$, so both product sequences are constant from that stage
onward. Their limit values retain all the inductive properties.

Otherwise those indices are cofinal. For each $\xi<\lambda$, choose
$i$ with $\xi\leq i<\lambda$ and $a_i>1$. Since $\lambda$ is a limit,
$i+1<\lambda$. The inductive bound and monotonicity give
\begin{equation}\label{eq:catch-up}
 J_\xi<O_\xi\cdot2\leq O_i\cdot2
         \leq O_i a_i=O_{i+1}\leq O_\lambda.
\end{equation}
Taking the supremum over $\xi$ yields $J_\lambda\leq O_\lambda$.
The reverse inequality follows from $O_\xi\leq J_\xi$ at earlier stages.
Hence the limit values agree. Their supports, leading monomials, and
coefficients therefore agree as well; the strict factor-two bound is
immediate for a positive ordinal compared with itself. This completes
the induction.
\end{proof}

Leading-monomial estimates also play a role in the natural-product
analysis in \cite[Lemmas~2.4--2.5]{lip-product}. Here the invariant
compares \emph{ordinary and Jacobsthal} partial products at arbitrary
ordinal stages.

The critical point is \eqref{eq:catch-up}. It does not presume that natural
sum or Jacobsthal multiplication is continuous in an inappropriate
argument. A later nonunit ordinary factor overtakes any current
Jacobsthal value, because the discrepancy never reaches the next doubling
of the leading scale.

\subsection{Why the cofinally active limit is a pure power}

\begin{proposition}\label{prop:pure-limit}
Under the cofinality hypothesis of \eqref{eq:limit-equality}, the common
value $p=O_\lambda=J_\lambda$ has the form $\omega^\delta$ with $\delta>0$.
\end{proposition}

\begin{proof}
Every partial product is strictly below $p$. Indeed, after any stage
there is a nonunit factor, which makes a subsequent ordinary partial
product strictly larger. Given $u,v<p$, select $\xi<\lambda$ with
$u,v<O_\xi$. Choose a later nonunit factor as in the preceding proof.
Then
\[
 u+v<O_\xi\cdot2\leq O_{i+1}<p.
\]
Thus $p$ is additively indecomposable. For completeness, the CNF criterion
for a positive ordinal to have this property is that it be a pure power
of $\omega$: if its leading coefficient is at least two, or if it has a
nonzero lower tail, one can split off a leading monomial to express it
as the sum of two smaller ordinals. Conversely, ordinals below a pure
power have sum below that power, by ordinary CNF addition.

Cofinitely many stages are not needed, but cofinally many nonunit stages
ensure that the value exceeds every finite integer: choose arbitrarily
many such stages and compare with the corresponding product of $2$'s.
Hence $p\geq\omega$, giving $\delta>0$.
\end{proof}

\begin{example}[A limit index by itself is insufficient]\label{ex:units}
Let $a_0=\omega+1$, $a_1=2$, and $a_i=1$ for $2\leq i<\omega$.
Then
\[
 O_\omega=\omega\cdot2+1,
 \qquad J_\omega=\omega\cdot2+2.
\]
Both sequences stabilize after the second factor. The ambient index is
a limit, but the active index is the finite ordinal $2$.
\end{example}

\begin{remark}
Finite degree sums do not directly extend to all transfinite stages.
For example, every factor in the constant sequence $2,2,\ldots$ has degree
zero, but its ordinary and Jacobsthal products at $\omega$ are both
$\omega$, of degree one. The colexicographic prefix construction avoids
any need for a false transfinite extension of \eqref{eq:metadata}.
\end{remark}

\section{Removing units and isolating the finite active tail}\label{sec:active}

\subsection{The active index}

\begin{definition}
For a nonzero sequence $(a_i)_{i<\kappa}$ let
\[
 S=\{i<\kappa:a_i>1\}.
\]
Give $S$ its inherited well-order, write $\tau=\otp(S)$, and let
$e:\tau\to S$ be its increasing enumeration. Put $b_\eta=a_{e(\eta)}$.
Write uniquely
\begin{equation}\label{eq:active-tail}
 \tau=\lambda+n,\qquad n<\omega,
\end{equation}
where $\lambda$ is zero or a nonzero limit. The factors indexed by
$\lambda,\ldots,\lambda+n-1$ form the \emph{finite active suffix}.
\end{definition}

The decomposition \eqref{eq:active-tail} is simply the removal of the
finite terminal part of the ordinal $\tau$. It is not the assertion
that all but finitely many original factors equal $1$. There may be
arbitrarily many nonunit factors in the initial part $\lambda$.

\begin{lemma}[Unit erasure]\label{lem:unit-erasure}
Deleting all factors equal to $1$ does not change either the ordinary
or the Jacobsthal product.
\end{lemma}

\begin{proof}
For $\xi\leq\kappa$, let $s(\xi)=\otp(S\cap\xi)$. Compare the original
partial product at $\xi$ with the active partial product at $s(\xi)$.
Induction on $\xi$ proves equality for each of the two operations.
At a successor, either the new factor is $1$ and neither active count
nor value changes, or the new active factor is appended on both sides.
At a limit, if the active indices below it are bounded, the values
stabilize after the last relevant initial segment. Otherwise their
counts are cofinal in the limit active count, and the two suprema are
taken over cofinal copies of the same nondecreasing sequence. This
includes the possibility of infinitely many active indices in a bounded
initial segment; after that segment the value is simply constant.
\end{proof}

\begin{theorem}[Finite-tail reduction]\label{thm:tail-reduction}
With the above notation, let
\[
 p=\prod_{\eta<\lambda}b_\eta
\]
be the ordinary product, with $p=1$ when $\lambda=0$. Then
\begin{align}\label{eq:tail-reduction}
 J((a_i)_{i<\kappa})
    &=p\,F(b_\lambda,\ldots,b_{\lambda+n-1}),\notag\\
 O((a_i)_{i<\kappa})
    &=p\,(b_\lambda\cdots b_{\lambda+n-1}).
\end{align}
When $\lambda>0$, $p=\omega^\delta$ for some $\delta>0$. Empty suffixes
have value $1$.
\end{theorem}

\begin{proof}
First erase units. If $\lambda=0$, the first identity is the finite
identification in Corollary~\ref{cor:finite-assoc}; the second is the
definition of a finite ordinary product.

Suppose $\lambda>0$. Every active factor is greater than $1$, so
Theorem~\ref{thm:sync} and Proposition~\ref{prop:pure-limit} give
$J_\lambda=O_\lambda=p=\omega^\delta$. Only $n$ successor stages remain.
Finite associativity, already proved from the cell construction, yields
\[
 J_{\lambda+n}
   =p\J F(b_\lambda,\ldots,b_{\lambda+n-1})
   =p\,F(b_\lambda,\ldots,b_{\lambda+n-1}).
\]
The last equality follows from Corollary~\ref{cor:binary-equality}, because
$p$ is a monomial. Ordinary associativity gives the second identity.
\end{proof}

\begin{corollary}\label{cor:equality-tail}
The full ordinary and Jacobsthal products are equal if and only if their
finite active suffixes have equal ordinary and Jacobsthal products.
\end{corollary}

\begin{proof}
Left multiplication by a fixed nonzero ordinal is strictly increasing
in the right operand. The common nonzero factor $p$ in
\eqref{eq:tail-reduction} can therefore be cancelled for equality.
\end{proof}

The reduction is stronger than a numerical estimate: all possible
Jacobsthal-specific behavior can be moved into one finite product at the
end of the active ordering. The prefix need not be small, countable, or
computably presented.

\section{A direct model for every ordinal-indexed product}\label{sec:transfinite}

\subsection{The ordinary finite-support prefix}

Let $(A_i)_{i<\theta}$ be nonempty well-orders, each with its least element
denoted $0_i$. Define
\begin{equation}\label{eq:fs-set}
 \fs((A_i)_{i<\theta})
 =\left\{f\in\prod_{i<\theta}A_i:
        \{i:f(i)\ne0_i\}\text{ is finite}\right\}.
\end{equation}
For distinct $f,g$, their finite union of supports contains a largest
coordinate at which they differ. Put $f<_{\CC}g$ when $f(i)<g(i)$ at
that coordinate. Write $\CC((A_i)_{i<\theta})$ for this ordered set.
The comparison rule is defined simultaneously on the full set
\eqref{eq:fs-set}.

\begin{lemma}[Standard finite-support model]\label{lem:ordinary-model}
The order $\CC((A_i)_{i<\theta})$ is a well-order extending the
componentwise partial order. Its type is the ordinary ordinal product
$\prod_{i<\theta}\otp(A_i)$.
\end{lemma}

\begin{proof}
The componentwise assertion follows immediately from the largest-difference
rule. For the type and well-order assertions, induct on $\theta$.
At zero the set is a singleton. At a successor, it is the ordinary colex
product of the previous finite-support order and the last factor, so its
type is multiplied on the right by that factor's type.

At a limit, every finite support is contained in some proper initial
segment of the index. The orders on these initial segments, extended by
zero on later coordinates, form an increasing family of \emph{initial
suborders}: any function with a nonzero later coordinate is larger than
all functions supported in the earlier segment. Their union is therefore
a well-order of type the supremum of their types. This proves the
recursion for the ordinary product. It is a proof of the direct model,
not its definition.
\end{proof}

This construction is recalled, with references to earlier sources, in
\cite[Section~3.2]{lip-product}. Finite support is essential. For example,
the full Cartesian product of countably many copies of $2$ has uncountable
cardinality, whereas the ordinary product $2^\omega$ is $\omega$.

\subsection{The final comparison rule}

Return to the original nonzero factors and their active enumeration.
Use the corresponding well-orders $B_\eta$ for $b_\eta$. Define
\[
 \BB=\CC((B_\eta)_{\eta<\lambda}),\qquad
 \mathfrak T=\FF(B_\lambda,\ldots,B_{\lambda+n-1}).
\]
When a list or prefix is empty, the associated order is a singleton.

\begin{definition}[The transfinite all-at-once order]\label{def:L}
On
$\fs((B_\eta)_{\eta<\tau})$, split a function into its prefix $f^-$
on $\lambda$ and its finite terminal tuple $f^+$. Put
\begin{equation}\label{eq:L-rule}
 f<_{\LL}g\quad\Longleftrightarrow\quad
 \bigl(f^+<_{\mathfrak T}g^+\bigr)
 \ \text{or}\
 \bigl(f^+=g^+\ \text{and}\ f^-<_{\BB}g^-\bigr).
\end{equation}
Restore the original unit coordinates, each of which has only its least
point, to regard this as an order on
$\fs((A_i)_{i<\kappa})$. Denote it by $\LL((A_i)_{i<\kappa})$.
\end{definition}

This is simply the ordinary colex product $\colex(\BB,\mathfrak T)$,
with the terminal tuple the more significant coordinate. No ordinal
value $J_\xi$, no iteration of $\J$, and no supremum occurs in the
comparison rule. Even the order type of $\BB$ need not be calculated to
compare two given prefix functions.

\begin{theorem}[All-at-once representation]\label{thm:transfinite}
For every set-sized ordinal-indexed family of nonempty well-orders,
$\LL$ is a well-order extending the componentwise order, and
\begin{equation}\label{eq:L-type}
 \otp\bigl(\LL((A_i)_{i<\kappa})\bigr)
     =J((\otp(A_i))_{i<\kappa}).
\end{equation}
The construction is invariant under order isomorphisms of the index
and of all factors.
\end{theorem}

\begin{proof}
By Lemma~\ref{lem:ordinary-model}, $\BB$ is a well-order of type $p$.
The finite order $\mathfrak T$ has type
$F(b_\lambda,\ldots,b_{\lambda+n-1})$. Their ordinary colex product is a
well-order of the product of these types. Theorem~\ref{thm:tail-reduction}
identifies this type with the desired Jacobsthal product.

If $f\leq g$ componentwise, their suffixes are componentwise ordered.
If the suffixes differ, Theorem~\ref{thm:extension} compares them in the
required direction. If they agree, Lemma~\ref{lem:ordinary-model} compares
the prefixes in that direction. Thus \eqref{eq:L-rule} extends the full
componentwise order. Restoring unit coordinates is a bijection preserving
that partial order.

Every operation in the definition is preserved by isomorphisms: the
active index, its finite terminal part, the canonical CNF blocks, their
exponents, the least points, finite support, and largest-coordinate
comparison. This proves invariance.
\end{proof}

\begin{corollary}[Finite-cell maximality in the transfinite case]
If the active suffix is nonempty, partition its tuple set into its finite
pivot cells $C$. Partition the full finite-support set into the
corresponding subsets $\BB\times C$, ordered internally by ordinary
colex. Among componentwise linear extensions having these specified
within-cell orders, $\LL$ has the largest order type. If the suffix is
empty, its one prescribed cell is the ordinary prefix order itself.
\end{corollary}

\begin{proof}
The type of $\BB$ is a pure power $\omega^\delta$, allowing $\delta=0$
when the prefix is empty. A suffix cell of type $\omega^e$ gives a full
cell of type $\omega^{\delta+e}$. Since $e\mapsto\delta+e$ is strictly
increasing, the decreasing-exponent order of the suffix cells also sorts
the full cells correctly. The same finite interleaving bound
(Lemma~\ref{lem:shuffle}) proves maximality.
\end{proof}

\subsection{Zero factors and the empty product}\label{subsec:zero}

The empty product is a singleton of type $1$. If some factor is empty,
the corresponding Cartesian set of finite-support functions is empty,
and we define $\LL$ to be empty, of type $0$. On the arithmetic side use
the usual absorbing-zero convention: any product with a zero factor is
zero. This agrees with successor multiplication and with the eventual
limit value after the zero has occurred.

One must not replace this convention by the supremum of \emph{all}
earlier partial products after a drop to zero. For instance, a sequence
beginning with a positive factor and then a zero has eventual value zero,
although the supremum of its entire earlier history is positive. All
supremum arguments in this manuscript were explicitly restricted to
nonzero factors, where the partial products are nondecreasing.

\subsection{What has, and has not, been made nonrecursive}

Definitions~\ref{def:F} and~\ref{def:L} are nonrecursive with respect to
Jacobsthal arithmetic. Their proofs of correctness use finite and
transfinite induction, just as the correctness proof for the ordinary
finite-support construction does. That is not a limit \emph{definition}
of the ordered set.

The finite regrouping theorem supplies a direct common-object proof of
ordinary binary associativity for $\J$. The transfinite model supplies a
single ordered set for every arity. It also inherits the known numerical
generalized associativity law once its type has been identified. We do
not assert that flattening an arbitrarily nested family is automatically
an order-preserving map, or give a separate pointwise formula for all
such transfinite regrouping isomorphisms.

\section{An exact three-state equality test}\label{sec:automaton}

\subsection{States and transitions}

At a nonzero partial stage compare the ordinary and Jacobsthal values.
Synchronization leaves only three relevant possibilities:
\[
 \mathsf M:\ \text{equal and monomial},\qquad
 \mathsf E:\ \text{equal and nonmonomial},\qquad
 \mathsf D:\ \text{different}.
\]
In state $\mathsf D$ neither value can be a monomial, because their
supports and leading coefficients agree. Let $q(b)$ be the number of
\emph{distinct positive exponents} in the CNF of a nonzero factor $b$,
and let $t=t(b)$.

\begin{theorem}[Transition table]\label{thm:automaton}
Starting from state $\mathsf M$ at the empty product, a finite list of
nonzero factors is classified exactly by Table~\ref{tab:states}.
For an arbitrary ordinal-indexed nonzero list, start from $\mathsf M$
and read only the finite active suffix. The full ordinary and Jacobsthal
products agree exactly when the final state is not $\mathsf D$.
\end{theorem}

\begin{table}[htbp]
\centering
\caption{Exact transitions after multiplication by a nonzero factor $b$.}
\label{tab:states}
\begin{tabular}{@{}p{0.48\textwidth}ccc@{}}
\toprule
Condition on $b$ & From $\mathsf M$ & From $\mathsf E$ & From $\mathsf D$\\
\midrule
$t=0$, $q(b)=1$ & $\mathsf M$ & $\mathsf M$ & $\mathsf M$\\
$t=0$, $q(b)\geq2$ & $\mathsf E$ & $\mathsf E$ & $\mathsf E$\\
$t=1$, $q(b)=0$ (the unit $1$) & $\mathsf M$ & $\mathsf E$ & $\mathsf D$\\
$t=1$, $q(b)\geq1$ & $\mathsf E$ & $\mathsf E$ & $\mathsf D$\\
$t\geq2$, $q(b)=0$ & $\mathsf M$ & $\mathsf D$ & $\mathsf D$\\
$t\geq2$, $q(b)\geq1$ & $\mathsf E$ & $\mathsf D$ & $\mathsf D$\\
\bottomrule
\end{tabular}
\end{table}

\begin{proof}
Let $x,y$ be the synchronized ordinary and Jacobsthal values before the
step. If $t=0$, Lemma~\ref{lem:one-step} shows that both new values are
$\omega^{d(x)}b$. Strict increase of $\sigma\mapsto d(x)+\sigma$
means their support has exactly $q(b)$ elements. Thus they reset to
$\mathsf M$ when $q(b)=1$ and to $\mathsf E$ otherwise.

Suppose $t>0$. The old exponent support survives in the terminal
contribution, and the $q(b)$ new positive-block exponents are distinct
and all lie strictly above it. Therefore the new support size is
\begin{equation}\label{eq:support-size}
 s_{\mathrm{new}}=s_{\mathrm{old}}+q(b).
\end{equation}
If $x\neq y$, their leading terms agree, so some lower coefficient of
$y$ is strictly greater than that of $x$. That inequality survives and
can only grow when all the coefficients of $y$ are multiplied by $t$.
Hence state $\mathsf D$ persists for $t>0$.

If $x=y$, the outputs are equal exactly when $t=1$ or $x$ is a monomial,
by the calculation behind Corollary~\ref{cor:binary-equality}. Combining
this condition with \eqref{eq:support-size} gives all four remaining
rows of the table. The initial value $1$ is monomial, so the initial
state is $\mathsf M$.

Finally, Corollary~\ref{cor:equality-tail} reduces equality of arbitrary
products to equality on the finite active suffix. The same finite table
therefore decides it. Equivalently, the preceding active limit produces
an equal pure-power prefix and thus resets the state to $\mathsf M$.
\end{proof}

\subsection{Interpretation and examples}

The automaton has two sources of reset. A nonzero limit factor removes
all dependence on previous lower coefficients. A cofinally active limit
of stages does the same, and additionally makes the result a pure power.
Between resets, discrepancy is created precisely when a finite tail of
size at least two multiplies an already nonmonomial equal value.
Once discrepancy exists, a factor with positive finite tail cannot
remove it.

For example, the active word $(\omega+1,2)$ gives
$\mathsf M\to\mathsf E\to\mathsf D$, so its two products differ.
Appending $\omega$ gives a reset to $\mathsf M$:
\[
 ((\omega+1)\cdot2)\omega=\omega^2
   =((\omega+1)\J2)\J\omega.
\]
Appending $\omega^2+\omega$ instead resets to $\mathsf E$: the new values
agree but have two positive-exponent terms. A later factor $2$ creates
discrepancy again.

Only the categories of $t(b)$ and $q(b)$ enter the equality decision.
The actual ordinal sizes of the exponents and the positive-exponent
coefficients do not. This does not mean that the product values are
independent of those data; it means precisely that the yes/no comparison
with the ordinary product is independent of them once the state is known.
The theorem is a finite decision procedure only when the finite active
suffix itself is available. It is not an algorithm for discovering that
suffix from an arbitrary oracle presentation of an infinite sequence.

\section{Worked families and boundary cases}\label{sec:examples}

\subsection{Three factors with visible cell contributions}

Let $a=\omega+2$. Its square illustrates the distinction between all
three multiplications:
\begin{align*}
 aa&=\omega^2+\omega\cdot2+2,\\
 a\J a&=\omega^2+\omega\cdot2+4,\\
 a\otimes a&=\omega^2+\omega\cdot4+4.
\end{align*}
For the direct triple product, a pivot at coordinate $3$ contributes
one cell of type $\omega^3$. A pivot at coordinate $2$ contributes two
cells of type $\omega^2$, one for each terminal point at coordinate $3$.
A pivot at coordinate $1$ contributes four $\omega$-cells and eight
singletons. Thus
\begin{equation}\label{eq:a-cube}
 F(a,a,a)=\omega^3+\omega^2\cdot2+\omega\cdot4+8.
\end{equation}
The ordinary cube is
$\omega^3+\omega^2\cdot2+\omega\cdot2+2$. The support and leading
monomial agree, while the lower coefficients differ exactly as the
synchronization theorem permits.

\subsection{A closed form for \texorpdfstring{$(\omega+c)^{\J n}$}{finite Jacobsthal powers of omega+c}}

For $c\geq1$ finite and $n\geq1$, write $J_n(c)$ for the product of $n$
copies of $\omega+c$. The cell formula immediately gives
\begin{equation}\label{eq:wc-finite}
 J_n(c)=\sum_{j=n}^{0}\omega^j c^{\,n-j}.
\end{equation}
In contrast,
\begin{align}\label{eq:wc-ordinary-natural}
 (\omega+c)^n
    &=\omega^n+\sum_{j=n-1}^{0}\omega^j c,\notag\\
 \underbrace{(\omega+c)\otimes\cdots\otimes(\omega+c)}_{n\text{ factors}}
    &=\sum_{j=n}^{0}\omega^j\binom{n}{j}c^{\,n-j}.
\end{align}
Every descending summation here means ordinary CNF concatenation in
the displayed decreasing order of exponents. Thus the formulas also
apply without ambiguous ellipses when $n=1$.

To verify \eqref{eq:wc-finite} directly, the final pivot contributes
$\omega^n$, the preceding pivot contributes $c$ copies at exponent
$n-1$, and so on. The finite-only tuples contribute $c^n$ singletons.
The ordinary formula follows by induction using
\eqref{eq:finite-ordinary}. The natural formula is the finite binomial
theorem for the natural semiring; its exponents here are finite, so
natural and ordinary addition of exponents agree.

These are applications of the construction, not claims that the
numerical finite-power identities were previously unknown. They also
show that the coefficient differences can be very large while
$O_n\leq J_n<O_n\cdot2$ still holds. An ordinal-scale bound is not a
coefficientwise ratio bound.

\subsection{Passing through a genuine limit and then taking successors}

For the constant sequence $a=\omega+2$ indexed by $\omega$, all indices
are active. The common prefix product is $\omega^\omega$. At index
$\omega+2$, the last two factors constitute the finite active suffix.
Hence
\begin{align*}
 J_{\omega+2}
 &=\omega^{\omega+2}+\omega^{\omega+1}\cdot2+\omega^\omega\cdot4,\\
 O_{\omega+2}
 &=\omega^{\omega+2}+\omega^{\omega+1}\cdot2+\omega^\omega\cdot2.
\end{align*}
Discrepancy has disappeared at $\omega$ and reappeared after two successor
factors. The automaton reads only $(\omega+2,\omega+2)$ and ends in
$\mathsf D$.

More generally, for any $a>1$ and $\beta=\lambda+n$ with $\lambda=0$
or a nonzero limit,
\begin{equation}\label{eq:power-reduction}
 a^{\J\beta}=a^\lambda
       F(\underbrace{a,\ldots,a}_{n\text{ copies}}).
\end{equation}
This is Jacobsthal exponentiation, the iteration of $\J$. It is not
super-Jacobsthal exponentiation, which iterates $\otimes$.
The direct finite-support construction therefore also supplies a model
for the operation in \eqref{eq:power-reduction}.

\subsection{Exponent absorption and the separation of pivot layers}

Suppose the degrees of successive factors include $\omega$ and $\omega$.
Their ordinary sum is $\omega\cdot2$, not a polynomial operation on
integer degrees. Suppose instead a prefix degree is finite and a later
pivot exponent is $\omega$; that finite prefix is absorbed:
$k+\omega=\omega$. Such absorption affects the calculated type and must
be handled correctly in any implementation of the cell formula.

It does \emph{not}, however, cause a collision between different pivot
layers. Lemma~\ref{lem:pivot-separation} shows that a pivot-$j$ exponent
lies above all preceding prefix degrees and at most $D_j$; all later
pivot layers lie strictly higher. Within one pivot layer, distinct input
block exponents remain distinct after addition of the common prefix
degree. Equal cell exponents therefore arise only from repeated blocks
of the same exponent and their different finite tail choices.

This separation explains why the primitive comparison rule needs no
prefix arithmetic at all. It also prevents a tempting but incorrect
concern that ordinary exponent absorption might reverse the order of
pivots. Regrouping still requires matching ordered blocks rather than
raw coordinate tags: the inner colex orders change under grouping, as
Remark~\ref{rem:flatten} demonstrates.

\subsection{Why an unrestricted infinite maximum is the wrong target}

The componentwise order on finite-support functions from $\omega$ to
$2$ has an infinite antichain: the functions supported at a single
coordinate. Thus the finite Cartesian-product maximal-order theorem
cannot simply be applied to the full finite-support product. Lipparini
discusses this obstruction and the need for additional restrictions in
\cite[Section~3.4]{lip-product}.

Our transfinite maximality statement deliberately fixes the ordinary
order within its finitely many large prefix--suffix cells. It does not
claim that Jacobsthal multiplication, or infinitary natural multiplication,
is the maximum over all linear extensions of the unrestricted componentwise
finite-support order.

\section{Reproducible checks and their limitations}\label{sec:checks}

\subsection{The exact ordinal kernel}

The archive contains a standard-library-only Python implementation of
hereditary Cantor normal forms below $\ep$. An ordinal is an immutable
finite tuple of strictly decreasing exponent--coefficient pairs; each
exponent is itself an ordinal in the same representation. Zero is the
empty tuple. The implementation provides ordinary addition and
multiplication, natural addition and multiplication, the classical binary
Jacobsthal formula, and a separate implementation of the all-arity pivot
formula.

The pivot routine never invokes the binary Jacobsthal routine. It computes
all ordinary prefix degrees, all finite suffix-tail products, and the
resulting cell monomials, then combines equal exponents. This provides a
useful comparison between independently organized formulas, although they
share the same low-level ordinal kernel.

The actual tuple comparison rule is also implemented for sampled points
in canonical blocks. A point records its block index, block exponent,
offset within the block, and position in the factor. This allows checks
of the componentwise-extension property independently of the type formula.
A second comparator uses the calculated cell degrees; it is checked
against the primitive, prefix-arithmetic-free rule on every tested pair.

\subsection{Test families and recorded results}

The test run in the archive used random seed $20260920$. All assertions
passed. Table~\ref{tab:tests} records the executed counts, rather than
planned counts.

\begin{table}[htbp]
\centering
\caption{Exact checks recorded in \texttt{results/verification\_results.json}.}
\label{tab:tests}
\begin{tabular}{@{}p{0.77\textwidth}r@{}}
\toprule
Test family & Count\\
\midrule
Binary pairs, including zero, over the small CNF domain & 729\\
Exhaustive triples over 26 positive polynomials below $\omega^3$ & 17,576\\
Additional finite lists with nested transfinite exponents & 3,500\\
Finite split/regrouping comparisons & 86,216\\
Explicit lists containing a zero factor & 15\\
Explicit empty-list case & 1\\
Sampled or exhaustively enumerated tuple-pair comparisons & 56,769\\
Agreements of primitive and cell-degree comparison rules & 56,769\\
Componentwise-comparable pairs within that set & 12,107\\
Additional tuple-triple transitivity checks & 4,000\\
\bottomrule
\end{tabular}
\end{table}

The small domain contains all nonzero expressions
$\omega^2c_2+\omega c_1+c_0$ with $c_i\in\{0,1,2\}$, giving $3^3-1=26$
values. The additional lists have lengths from zero to seven and use
exponents selected from a pool containing $0,1,2,3,\omega,\omega+1,
\omega\cdot2,\omega^2,\omega^\omega$, and $\omega^\omega+\omega$.
Positive coefficients range from one to four.

For every nonzero list, the test checks the direct formula against binary
iteration; all finite split points; equality of supports and leading
monomials; coefficient domination; the strict factor-two bound; the degree
and finite-tail formulas; the natural-product upper bound; and the
three-state classification. Binary tests also check the defining successor
equation and the exact binary equality criterion.

The tuple tests check comparison antisymmetry, separation of distinct
points, componentwise extension, and sampled transitivity. Some small
families are exhausted, but only among the selected points; an infinite
block is never replaced by a claim that its finite sample is exhaustive.

\subsection{What the tests cannot establish}

The computational domain is below $\ep$, whereas the theorems are stated
for all ordinals. No arbitrary transfinite limit is evaluated by the
program. The displayed $\omega+2$ examples are computed \emph{after}
substituting the prefix value justified by the proof. In particular,
random tests do not certify the limit argument, the pure-power theorem,
or the correctness of the reduction for arbitrary cofinalities.

The implementation is not a formalization in Lean, Isabelle, or another
proof assistant. Sharing the low-level CNF routines creates a possibility
of correlated implementation mistakes. The checks should be viewed as
structured attempts to falsify finite consequences and catch transcription
errors, not as substitutes for the English proofs.

\section{Assessment, scope, and remaining questions}\label{sec:scope}

\subsection{The precise positive result}

The manuscript gives an affirmative construction for the following
precise version of the selected question: there is an isomorphism-invariant,
CNF-based well-order on finite tuples, and on finite-support tuples at
arbitrary ordinal arities, defined without Jacobsthal iteration or its
limit value, whose type is the Jacobsthal product. The finite objects
prove associativity through an explicit comparison of their blocks.
The restricted maximality theorems describe the same values without
choosing a tie-breaking order.

The binary comparison with natural multiplication also has the requested
order-theoretic explanation: the constructed order is one componentwise
linear extension, while the natural product is the maximum among all
such extensions. This addresses the binary inequality part of
Question~5.2, not every question listed there.

These statements are stronger than exhibiting a CNF formula for a
parenthesized product: the underlying sets, cells, comparison rules,
componentwise compatibility, and regrouping argument have all been
specified. On the other hand, the adjective ``natural'' in the original
question has no unique formal definition. The construction uses canonical
Cantor blocks and should be judged against that explicit design choice.

\subsection{Claims not made}

We do not claim the classical associativity law, the classical binary
CNF formula, the ordinary finite-support representation, or the numerical
inequalities as new. We do not claim that our comparison rule is preserved
by arbitrary embeddings, that it is independent of canonical CNF
structure, or that the tuple-flattening map is the associator.
We do not settle Altman's separate questions about super-Jacobsthal
exponentiation or all infinitary natural-product inequalities.

The literature search did not establish a priority theorem. In
particular, the synchronization and finite-tail observations may have
antecedents not found by the targeted search. The contribution presented
for evaluation is the combination of the concrete pivot model,
its finite-cell extremal description, its blockwise associativity proof,
and the arbitrary-index finite-tail representation, with all dependencies
made explicit.

\subsection{Further research directions suggested by the construction}

A useful next mathematical question is whether the pivot-cell constraints
can be replaced by an intrinsic condition on order embeddings, convex
suborders, or initial segments, without first displaying the CNF blocks.
Such a characterization would clarify how canonical this model is among
other possible answers to the same informal question.

A second question is how to describe the unique regrouping isomorphisms
by explicit coordinate transformations. The example in
Remark~\ref{rem:flatten} shows that coordinate flattening is insufficient;
some absorption isomorphisms inside monomial cells are unavoidable.
An effective description for hereditary CNFs could be suitable for
formalization.

A third direction is to identify the exact relationship between the
finite-suffix maximality condition here and the finite-prefix conditions
used for countable natural products. Their orientations differ: the
Jacobsthal discrepancy is confined to the \emph{end} of the active index,
whereas the countable natural-product construction in
\cite{lip-product} uses a finite initial part followed by ordinary
infinite behavior. This comparison is a research direction, not a claimed
new equivalence.

\section{Conclusion}

The main obstruction to a direct infinitary Jacobsthal model turns out,
in this approach, not to be an infinite combinatorial object of a new
kind. The relevant infinite object is the ordinary finite-support product.
What is different is a finite terminal rearrangement, organized by the
last coordinate occupying a positive-exponent Cantor block.

The finite partition makes the all-arity object explicit. The common
leading monomial bounds every discrepancy by the next ordinary doubling.
Cofinal activity then forces the two operations to agree at a limit,
leaving only finitely many active successor factors to be handled by the
pivot construction. Together these observations yield a concrete proposed
answer to the selected open-ended order-theoretic question, with a clear
account of both its mathematical strength and its limitations.

\appendix
\section{Proof dependency map and audit checklist}\label{app:audit}

The mathematical dependency chain is as follows. Ordinary CNF absorption
proves the classical binary formula directly from the defining recursion.
The pivot partition and ordinary cell orders give their monomial types.
The explicit finite cell formula then proves the binary identification,
the metadata lemma, and the finite regrouping theorem. None of these
steps invokes Jacobsthal associativity as a premise.

Separately, the binary formula gives the one-step synchronization lemma.
That lemma and the factor-two estimate are proved simultaneously through
all ordinal stages. At a cofinally active limit the estimate yields
mutual cofinal domination, hence equality; additive indecomposability then
yields a pure power. Unit erasure and finite associativity produce the
finite-tail reduction. The direct ordinary finite-support model, together
with the finite pivot model, gives the final all-at-once order.

The three most delicate checks are worth isolating. First, the global
pivot can only move left under a coordinatewise increase, and such a move
strictly lowers its cell exponent. Second, at a cofinally active limit,
for each earlier stage $\xi$ there really is a later successor stage
$i+1<\lambda$ that overtakes $J_\xi$; the limit hypothesis ensures this
successor is still below the endpoint. Third, the active index, not the
ambient ordinal index, controls the finite suffix. Example~\ref{ex:units}
shows exactly what would go wrong if units were not removed.

The mathematical argument is not restricted to countable ordinals.
Every invocation of cofinality chooses a later active index individually;
no countable cofinal sequence is assumed. The CNF decomposition has
finitely many terms for each input ordinal, regardless of its cardinality.
The number of factors may be any set-sized ordinal, but proper-class
products are outside the scope of the construction.

\section{Using and rebuilding the archive}\label{app:archive}

The principal files are \texttt{article.tex}, \texttt{article.pdf},
\texttt{code/ordinals.py}, and \texttt{code/verify.py}. The results directory
contains machine-readable and plain-text records of the executed checks.
The notes directory contains a short proof audit and a source-search log.
No downloaded third-party article PDFs or font files are redistributed.

From the archive's top-level directory, run
\begin{lstlisting}
python3 code/verify.py
latexmk -pdf -interaction=nonstopmode -halt-on-error article.tex
\end{lstlisting}
The supplied \texttt{Makefile} provides \texttt{make check},
\texttt{make pdf}, and \texttt{make clean}. The source has an internal
bibliography, so no separate bibliography processor is required. The
Python checks require only the standard library. The LaTeX document uses
common packages included in standard TeX Live installations.

\clearpage
\phantomsection
\addcontentsline{toc}{section}{References}
\begin{thebibliography}{9}

\bibitem{altman}
H. Altman,
\emph{Intermediate arithmetic operations on ordinal numbers},
Mathematical Logic Quarterly \textbf{63} (2017), 228--242.
DOI: \href{https://doi.org/10.1002/malq.201600006}{10.1002/malq.201600006}.
Author version: \href{https://arxiv.org/abs/1501.05747}{arXiv:1501.05747},
version 9, 18 September 2017. Questions 2.16 and 5.2 are the selected
problem references.

\bibitem{jacobsthal}
E. Jacobsthal,
\emph{Zur Arithmetik der transfiniten Zahlen},
Mathematische Annalen \textbf{67} (1909), 130--143.
The binary formula and algebraic laws used for historical context are
also restated in \cite{altman}; the binary formula is reproved here.

\bibitem{carruth}
P. W. Carruth,
\emph{Arithmetic of ordinals with applications to the theory of ordered
Abelian groups},
Bulletin of the American Mathematical Society \textbf{48} (1942), 262--271.
\href{https://projecteuclid.org/journals/bulletin-of-the-american-mathematical-society/volume-48/issue-4/Arithmetic-of-ordinals-with-applications-to-the-theory-of-ordered/bams/1183504256.full}{Project Euclid record}.
The finite-product maximum theorem is also explicitly recalled in
\cite[Section~3.3]{lip-product}.

\bibitem{lip-sums}
P. Lipparini,
\emph{Some transfinite natural sums},
Mathematical Logic Quarterly \textbf{64} (2018), 514--528.
\href{https://arxiv.org/abs/1509.04078}{arXiv:1509.04078}.
The earlier title cited by Altman was \emph{Some iterated natural sums}.

\bibitem{lip-product}
P. Lipparini,
\emph{An infinite natural product},
Annales Mathematicae Silesianae \textbf{32}(1) (2018), 247--262.
\href{https://arxiv.org/abs/1703.06908}{arXiv:1703.06908}, version 2.
Sections 3.2--3.4 describe the ordinary finite-support product, the
finite natural-product maximum theorem, and the obstruction to an
unrestricted infinitary extension.

\bibitem{lip-semigroup}
P. Lipparini,
\emph{Noncommutative infinitary semigroups},
\href{https://arxiv.org/abs/2605.28413}{arXiv:2605.28413},
version 1, 27 May 2026.
Cited as related general context, not as a source of the Jacobsthal
construction proposed here.

\end{thebibliography}
\end{document}
